% =====================================================================
\chapter{Soluciones de los ejercicios}
\label{ap:soluciones}
% =====================================================================

Cada solución se presenta en cinco partes: datos, planteamiento, desarrollo, resultado y comprobación. En las demostraciones, los datos son las hipótesis, el planteamiento es la estrategia y la comprobación es un caso particular o una verificación independiente. Los valores aproximados se dan con cuatro cifras significativas salvo que se indique otra cosa.

% ---------------------------------------------------------------------
\section{Tema 1}

\begin{solucion}{eje:1-clasificar}
\datos los números $-3$, $0{,}75$, $0{,}\overline3$, $\sqrt9$, $\sqrt8$, $\pi$, $1/7$, $0$.
\planteamiento escribir cada número como fracción si es posible (\cref{prop:decimal-periodico}).
\pasos
\begin{enumerate}[label=\arabic*.]
  \item $-3\in\Z$, no es natural. $0\in\Z$ y no está en $\N$ con el convenio de estos apuntes.
  \item $0{,}75=3/4$ y $0{,}\overline3=1/3$ son racionales no enteros; $1/7\in\Q$.
  \item $\sqrt9=3\in\N$.
  \item $\sqrt8=2\sqrt2$ es irracional: si fuera racional, $\sqrt2=\sqrt8/2$ también lo sería, en contra del \cref{teo:raiz2}. $\pi$ es irracional (resultado conocido, demostrado en~\cite[cap.~16]{spivak}).
\end{enumerate}
\resultado \textbf{$\N$: $\sqrt9$. $\Z$: además $-3$ y $0$. $\Q$: además $0{,}75$, $0{,}\overline3$ y $1/7$. $\R\setminus\Q$: $\sqrt8$ y $\pi$.} Todos son reales.
\comprobacion las inclusiones $\N\subset\Z\subset\Q$ hacen que cada número de un conjunto pertenezca también a los siguientes.
\end{solucion}

\begin{solucion}{eje:1-axiomas}
\datos $a+x=a+y$.
\planteamiento identificar en cada paso qué propiedad se usa.
\pasos
\begin{enumerate}[label=\arabic*.]
  \item (1) neutro de la suma (C3), con la conmutatividad (C2) para escribir $0+x=x+0$.
  \item (2) opuesto (C4), con (C2): $(-a)+a=a+(-a)=0$.
  \item (3) asociatividad (C1).
  \item (4) la hipótesis $a+x=a+y$ (no es un axioma).
  \item (5) asociatividad (C1). (6) opuesto (C4). (7) neutro (C3).
\end{enumerate}
\resultado \textbf{(1) C3; (2) C4; (3) C1; (4) hipótesis; (5) C1; (6) C4; (7) C3} (con C2 donde se cambia el orden).
\comprobacion es la ley de cancelación de la suma; con $y=0$ se obtiene la \cref{prop:cuerpo}(i).
\end{solucion}

\begin{solucion}{eje:1-vf}
\planteamiento buscar un contraejemplo o una demostración por reducción al absurdo.
\pasos
\begin{enumerate}[label=(\alph*)]
  \item Falsa: $\sqrt2$ y $-\sqrt2$ son irracionales y su suma es $0\in\Q$.
  \item Verdadera: si $r\neq0$ es racional, $x$ irracional y $rx=s\in\Q$, entonces $x=s/r\in\Q$, contradicción.
  \item Verdadera: si $a<b$ son racionales, $(a+b)/2$ es racional y está entre ambos (\cref{eje:1-epsilon}).
\end{enumerate}
\resultado \textbf{(a) falsa; (b) verdadera; (c) verdadera.}
\comprobacion en (c), con $a=\frac13$, $b=\frac12$: $\frac{5}{12}$ cumple $\frac4{12}<\frac5{12}<\frac6{12}$.
\end{solucion}

\begin{solucion}{eje:1-fracciones}
\datos $0{,}1\overline6$, $2{,}\overline{45}$, $0{,}\overline{142857}$.
\planteamiento multiplicar por potencias de~10 para eliminar el periodo (\cref{ej:periodico}).
\pasos
\begin{enumerate}[label=(\alph*)]
  \item $x=0{,}1\overline6$: $100x=16{,}\overline6$ y $10x=1{,}\overline6$; restando, $90x=15$, luego $x=15/90=1/6$.
  \item $x=2{,}\overline{45}$: $100x-x=245{,}\overline{45}-2{,}\overline{45}=243$, luego $x=243/99=27/11$.
  \item $x=0{,}\overline{142857}$: $10^6x-x=142857$, luego $x=142857/999999=1/7$ (dividiendo numerador y denominador por $142857$).
\end{enumerate}
\resultado \textbf{(a) $1/6$; (b) $27/11$; (c) $1/7$.}
\comprobacion $1\div6=0{,}1666\ldots$, $27\div11=2{,}4545\ldots$ y $1\div7=0{,}142857\,142857\ldots$
\end{solucion}

\begin{solucion}{eje:1-potencias}
\planteamiento escribir todo como potencias de primos y aplicar la \cref{prop:exponentes}.
\pasos
\begin{enumerate}[label=(\alph*)]
  \item $\dfrac{2^{-3}\cdot2^{4}}{2^{-3}}=2^{-3+4+3}=2^4=16$.
  \item $27^{2/3}=(3^3)^{2/3}=3^2=9$; $16^{-3/4}=(2^4)^{-3/4}=2^{-3}=\frac18$. Producto: $\frac98$.
  \item $\sqrt{50}=5\sqrt2$, $\sqrt{18}=3\sqrt2$, $\sqrt8=2\sqrt2$; total $(5-3+2)\sqrt2=4\sqrt2$.
\end{enumerate}
\resultado \textbf{(a) $16$; (b) $9/8$; (c) $4\sqrt2$.}
\comprobacion (c): $7{,}0711-4{,}2426+2{,}8284=5{,}6569\approx4\cdot1{,}4142$.
\end{solucion}

\begin{solucion}{eje:1-racionalizar}
\planteamiento multiplicar por el conjugado, o por el factor que completa una potencia exacta.
\pasos
\begin{enumerate}[label=(\alph*)]
  \item $\dfrac{1}{\sqrt3-\sqrt2}\cdot\dfrac{\sqrt3+\sqrt2}{\sqrt3+\sqrt2}=\dfrac{\sqrt3+\sqrt2}{3-2}=\sqrt3+\sqrt2$.
  \item $\dfrac{2(\sqrt5-1)}{(\sqrt5+1)(\sqrt5-1)}=\dfrac{2(\sqrt5-1)}{4}=\dfrac{\sqrt5-1}2$.
  \item $\dfrac{6}{2^{1/3}}\cdot\dfrac{2^{2/3}}{2^{2/3}}=\dfrac{6\cdot2^{2/3}}{2}=3\sqrt[3]4$.
\end{enumerate}
\resultado \textbf{(a) $\sqrt3+\sqrt2$; (b) $(\sqrt5-1)/2$; (c) $3\sqrt[3]4$.}
\comprobacion (b): $2/(1+2{,}2361)=0{,}6180$ y $(2{,}2361-1)/2=0{,}6180$.
\end{solucion}

\begin{solucion}{eje:1-ordenar}
\datos $3{,}1415<\pi<3{,}1416$.
\planteamiento los tres números son positivos; por el \cref{eje:1-potencia-n}, basta comparar sus cuadrados.
\pasos
\begin{enumerate}[label=\arabic*.]
  \item $(\sqrt2+\sqrt3)^2=5+2\sqrt6$ y $(\sqrt{10})^2=10$. Como $\sqrt6<2{,}5$ (porque $6<6{,}25$), $5+2\sqrt6<10$, luego $\sqrt2+\sqrt3<\sqrt{10}$.
  \item $\pi^2<3{,}1416^2=9{,}86965\ldots<9{,}8697$. Por otro lado, $\sqrt6>2{,}449$ porque $2{,}449^2=5{,}997601<6$, luego $5+2\sqrt6>9{,}898>9{,}8697>\pi^2$, y $\pi<\sqrt2+\sqrt3$.
\end{enumerate}
\resultado \textbf{$\pi<\sqrt2+\sqrt3<\sqrt{10}$.}
\comprobacion con calculadora: $3{,}1416<3{,}1463<3{,}1623$.
\end{solucion}

\begin{solucion}{eje:1-cientifica}
\planteamiento agrupar mantisas y potencias de~10 por separado.
\pasos
\begin{enumerate}[label=\arabic*.]
  \item Mantisas: $3{,}2\cdot4/1{,}6=12{,}8/1{,}6=8$. Potencias: $10^{5-8-(-2)}=10^{-1}$.
  \item (a) $8\cdot10^{-1}$.
  \item (b) $7\cdot10^{10}\cdot2=1{,}4\cdot10^{11}$ bytes $=1{,}4\cdot10^{11}/10^9=140$ GB.
\end{enumerate}
\resultado \textbf{(a) $8\cdot10^{-1}=0{,}8$; (b) $1{,}4\cdot10^{11}$ bytes, es decir, $140$ GB.}
\comprobacion (a): $3{,}2\cdot10^5\cdot4\cdot10^{-8}=1{,}28\cdot10^{-2}$, y $1{,}28\cdot10^{-2}/1{,}6\cdot10^{-2}=0{,}8$.
\end{solucion}

\begin{solucion}{eje:1-intervalos}
\datos $A=[-2,3)$, $B=(1,5]$.
\planteamiento dibujar ambos intervalos en la recta y respetar qué extremos se incluyen.
\pasos
\begin{enumerate}[label=\arabic*.]
  \item $A\cap B$: $x>1$ y $x<3$; $(1,3)$.
  \item $A\cup B$: desde $-2$ (incluido) hasta $5$ (incluido), sin huecos porque los intervalos se solapan; $[-2,5]$.
  \item $A\setminus B$: puntos de $A$ con $x\le1$; $[-2,1]$.
  \item $\R\setminus A$: $x<-2$ o $x\ge3$; $(-\infty,-2)\cup[3,\infty)$.
\end{enumerate}
\resultado \textbf{$(1,3)$; $[-2,5]$; $[-2,1]$; $(-\infty,-2)\cup[3,\infty)$.}
\comprobacion $1\in A\setminus B$ porque $1\in A$ y $1\notin B$; $3\notin A$, luego $3\in\R\setminus A$.
\end{solucion}

\begin{solucion}{eje:1-ecuacion-producto}
\datos $x^3=x$.
\planteamiento pasar todo a un miembro, factorizar y aplicar la \cref{prop:cuerpo}(iv).
\pasos
\begin{enumerate}[label=\arabic*.]
  \item $x^3-x=0$, es decir, $x(x^2-1)=x(x-1)(x+1)=0$.
  \item Por la \cref{prop:cuerpo}(iv), aplicada dos veces, algún factor es cero: $x=0$, $x=1$ o $x=-1$.
  \item Si se divide por $x$ se obtiene $x^2=1$, que pierde la solución $x=0$: la división solo es válida si $x\neq0$.
\end{enumerate}
\resultado \textbf{$x\in\{-1,0,1\}$.} Dividir por $x$ pierde $x=0$.
\comprobacion $(-1)^3=-1$, $0^3=0$, $1^3=1$.
\end{solucion}

\begin{solucion}{eje:1-quebrados}
\datos $b,d\neq0$; $a/b=ab^{-1}$; $(bd)^{-1}=b^{-1}d^{-1}$.
\planteamiento escribir cada cociente como producto por el inverso y usar (C1), (C2), (C6).
\pasos
\begin{enumerate}[label=\arabic*.]
  \item $bd\neq0$ por la \cref{prop:cuerpo}(iv), así que $(bd)^{-1}$ existe.
  \item $\frac{ad+bc}{bd}=(ad+bc)b^{-1}d^{-1}=adb^{-1}d^{-1}+bcb^{-1}d^{-1}$ (C6).
  \item Reordenando (C1, C2): $adb^{-1}d^{-1}=ab^{-1}(dd^{-1})=ab^{-1}$ y $bcb^{-1}d^{-1}=cd^{-1}(bb^{-1})=cd^{-1}$, por (C5) y (C3).
  \item Luego $\frac{ad+bc}{bd}=ab^{-1}+cd^{-1}=\frac ab+\frac cd$.
  \item Producto: $\frac ab\cdot\frac cd=ab^{-1}cd^{-1}=(ac)(b^{-1}d^{-1})=(ac)(bd)^{-1}=\frac{ac}{bd}$.
\end{enumerate}
\resultado \textbf{Ambas fórmulas se deducen de los axiomas de cuerpo.}
\comprobacion $\frac12+\frac13=\frac{1\cdot3+2\cdot1}{6}=\frac56$ y $\frac12\cdot\frac23=\frac26=\frac13$.
\end{solucion}

\begin{solucion}{eje:1-negativo}
\datos axiomas (O1)--(O3) y la \cref{prop:orden}.
\planteamiento reducir cada desigualdad a la pertenencia de una diferencia a $P$.
\pasos
\begin{enumerate}[label=(\alph*)]
  \item $c<0$ significa $-c\in P$, y $a<b$ significa $b-a\in P$. Por (O3), $(b-a)(-c)\in P$, es decir, $ac-bc\in P$, que es $bc<ac$.
  \item Sea $a>0$. Como $aa^{-1}=1\neq 0$, $a^{-1}\neq0$. Si fuera $a^{-1}<0$, por (a) con $0<a$ y $c=a^{-1}$ se tendría $0\cdot a^{-1}>a\cdot a^{-1}$, es decir, $0>1$, en contra de la \cref{prop:orden}(iv). Por tricotomía, $a^{-1}>0$.
  \item Si $0<a<b$, entonces $ab>0$ por (O3) y $(ab)^{-1}>0$ por (b). Multiplicando $a<b$ por $(ab)^{-1}$ (\cref{prop:orden}(iii)): $a(ab)^{-1}<b(ab)^{-1}$, es decir, $b^{-1}<a^{-1}$.
\end{enumerate}
\resultado \textbf{Las tres propiedades se deducen de los axiomas de orden.}
\comprobacion con $a=2$, $b=4$: $\frac14<\frac12$; con $c=-1$, $2<4$ da $-2>-4$.
\end{solucion}

\begin{solucion}{eje:1-potencia-n}
\datos $0\le a<b$.
\planteamiento inducción sobre $n$, multiplicando desigualdades de números no negativos.
\pasos
\begin{enumerate}[label=\arabic*.]
  \item Caso $n=1$: es la hipótesis.
  \item Paso: supóngase $a^n<b^n$. Como $a\ge0$, $a^{n+1}=a^n\cdot a\le a^n\cdot b$ (si $a=0$ es una igualdad $0=0\le\dots$; si $a>0$, se usa $a<b$ y $a^n\ge0$).
  \item Como $b>0$, $a^nb<b^nb=b^{n+1}$ por la \cref{prop:orden}(iii).
  \item Luego $a^{n+1}\le a^nb<b^{n+1}$.
  \item Recíproco para $n=2$: si $a,b\ge0$ y $a^2<b^2$, no puede ser $b\le a$, porque entonces $b^2\le a^2$ (por lo anterior si $b<a$, o igualdad si $b=a$). Por tricotomía, $a<b$.
\end{enumerate}
\resultado \textbf{$a^n<b^n$ para todo $n\in\N$; y para no negativos, $a^2<b^2\Rightarrow a<b$.}
\comprobacion la hipótesis $a\ge0$ es necesaria: $-3<2$ pero $(-3)^2>2^2$.
\end{solucion}

\begin{solucion}{eje:1-raiz3}
\datos se quiere probar que $\sqrt3\notin\Q$ y $\sqrt2+\sqrt3\notin\Q$.
\planteamiento reducción al absurdo, como en el \cref{teo:raiz2}, con el análogo del \cref{lem:par} para el~3.
\pasos
\begin{enumerate}[label=\arabic*.]
  \item Lema: si $3\nmid n$, entonces $n=3k\pm1$ y $n^2=9k^2\pm6k+1=3(3k^2\pm2k)+1$, que no es múltiplo de~3. Por contrarrecíproco, si $3\mid n^2$, entonces $3\mid n$.
  \item Si $\sqrt3=p/q$ irreducible, $p^2=3q^2$, luego $3\mid p$; con $p=3r$, $9r^2=3q^2$, $q^2=3r^2$ y $3\mid q$. Contradicción.
  \item Si $\sqrt2+\sqrt3=s\in\Q$, entonces $s\neq0$ y, elevando al cuadrado, $s^2=5+2\sqrt6$, luego $\sqrt6=(s^2-5)/2\in\Q$.
  \item Pero $\sqrt6$ es irracional: si $\sqrt6=p/q$ irreducible, $p^2=6q^2$ es par, $p=2r$, $2r^2=3q^2$, luego $q^2$ es par (porque $3q^2$ es par y 3 es impar) y $q$ es par. Contradicción.
\end{enumerate}
\resultado \textbf{$\sqrt3$ y $\sqrt2+\sqrt3$ son irracionales.}
\comprobacion el argumento falla, como debe, para $\sqrt4$: $p^2=4q^2$ no conduce a contradicción porque $p/q=2/1$ es irreducible con $q$ impar.
\end{solucion}

\begin{solucion}{eje:1-epsilon}
\datos (a) $a<b$; (b) $x\le y+\varepsilon$ para todo $\varepsilon>0$.
\planteamiento (a) operar con la \cref{prop:orden}; (b) reducción al absurdo eligiendo un $\varepsilon$ concreto.
\pasos
\begin{enumerate}[label=\arabic*.]
  \item (a) Sumando $a$ a $a<b$: $2a<a+b$. Multiplicando por $\frac12>0$ (positivo por el \cref{eje:1-negativo}(b), ya que $2=1+1>0$): $a<\frac{a+b}2$. Análogamente, sumando $b$: $a+b<2b$ y $\frac{a+b}2<b$.
  \item (b) Supóngase $x>y$. Entonces $\varepsilon=\frac{x-y}2>0$ y, por hipótesis, $x\le y+\frac{x-y}2=\frac{x+y}2$.
  \item Pero por (a) con $y<x$, $\frac{x+y}2<x$. Contradicción; luego $x\le y$.
\end{enumerate}
\resultado \textbf{(a) el punto medio está estrictamente entre $a$ y $b$; (b) $x\le y$.}
\comprobacion en (b), la conclusión no puede mejorarse a $x<y$: con $x=y$ la hipótesis se cumple.
\end{solucion}
% ---------------------------------------------------------------------
\section{Tema 2}

\begin{solucion}{eje:2-calcular}
\planteamiento aplicar la \cref{def:abs} según el signo de lo que hay dentro de las barras.
\pasos
\begin{enumerate}[label=\arabic*.]
  \item $\abs{-3}=3$ y $\abs{2-7}=\abs{-5}=5$; la diferencia es $-2$.
  \item $\pi-4<0$, luego $\abs{\pi-4}=4-\pi$.
  \item Si $x<1$, $x-1<0$ y $\abs{x-1}=1-x$.
\end{enumerate}
\resultado \textbf{$-2$; $4-\pi$; $1-x$.}
\comprobacion $4-\pi\approx0{,}858>0$, como debe ser un valor absoluto.
\end{solucion}

\begin{solucion}{eje:2-vf}
\planteamiento contrastar con el \cref{teo:triangular}, su \cref{cor:triangular-inversa} y la \cref{prop:abs}.
\pasos
\begin{enumerate}[label=(\alph*)]
  \item Falsa: con $a=0$, $b=1$, $\abs{a-b}=1>-1=\abs a-\abs b$. Lo correcto es $\abs{a-b}\ge\abs a-\abs b$.
  \item Verdadera: es el caso de igualdad del \cref{teo:triangular}.
  \item Falsa para $x<0$: $\sqrt{(-2)^2}=2\neq-2$. Lo correcto es $\sqrt{x^2}=\abs x$.
\end{enumerate}
\resultado \textbf{(a) falsa; (b) verdadera; (c) falsa.}
\comprobacion (b): $\abs{3+5}=8=3+5$ con $3\cdot5>0$; $\abs{3-5}=2\neq8$ con $3\cdot(-5)<0$.
\end{solucion}

\begin{solucion}{eje:2-traducir}
\planteamiento un intervalo $(c-r,c+r)$ es $\abs{x-c}<r$, con $c$ el punto medio y $r$ el radio.
\pasos
\begin{enumerate}[label=(\alph*)]
  \item Centro $(2+8)/2=5$, radio $3$: $\abs{x-5}<3$.
  \item Centro $(-1+4)/2=1{,}5$, radio $2{,}5$, extremos incluidos: $\abs{x-1{,}5}\le2{,}5$.
  \item $\abs{x-3}<0{,}01$.
\end{enumerate}
\resultado \textbf{(a) $\abs{x-5}<3$; (b) $\abs{x-1{,}5}\le2{,}5$; (c) $\abs{x-3}<0{,}01$.}
\comprobacion (b): $\abs{-1-1{,}5}=2{,}5$ y $\abs{4-1{,}5}=2{,}5$, los extremos cumplen la igualdad.
\end{solucion}

\begin{solucion}{eje:2-inec-simple}
\planteamiento (a) \cref{prop:abs}(iv); (b) $\abs A=\abs B$ si y solo si $A=B$ o $A=-B$.
\pasos
\begin{enumerate}[label=\arabic*.]
  \item (a) $-7<3x-2<7$; sumando 2: $-5<3x<9$; dividiendo por $3>0$: $-\frac53<x<3$.
  \item (b) Caso $2x-1=x+3$: $x=4$.
  \item Caso $2x-1=-(x+3)$: $3x=-2$, $x=-\frac23$.
\end{enumerate}
\resultado \textbf{(a) $\bigl(-\frac53,3\bigr)$; (b) $x=4$ o $x=-\frac23$.}
\comprobacion (b): $\abs{7}=\abs{7}$ y $\abs{-\frac73}=\abs{\frac73}$.
\end{solucion}

\begin{solucion}{eje:2-inec-casos}
\datos $\abs{x+1}\ge2x-4$.
\planteamiento separar los casos $x\ge-1$ y $x<-1$.
\pasos
\begin{enumerate}[label=\arabic*.]
  \item Si $x\ge-1$: $x+1\ge2x-4$, es decir, $x\le5$. Solución en este caso: $[-1,5]$.
  \item Si $x<-1$: $-x-1\ge2x-4$, es decir, $3\ge3x$, $x\le1$, que se cumple siempre en este caso. Solución: $(-\infty,-1)$.
  \item Unión: $(-\infty,5]$.
\end{enumerate}
\resultado \textbf{$x\in(-\infty,5]$.}
\comprobacion $x=5$: $6\ge6$; $x=6$: $7\ge8$ es falso; $x=-10$: $9\ge-24$.
\end{solucion}

\begin{solucion}{eje:2-inec-suma}
\datos $\abs{x-1}+\abs{x+2}\le5$; puntos de corte $-2$ y $1$.
\planteamiento estudiar los tres intervalos como en el \cref{ej:abs-casos}.
\pasos
\begin{enumerate}[label=\arabic*.]
  \item $x\ge1$: $2x+1\le5$, $x\le2$; da $[1,2]$.
  \item $-2\le x<1$: $3\le5$, siempre; da $[-2,1)$.
  \item $x<-2$: $-2x-1\le5$, $x\ge-3$; da $[-3,-2)$.
  \item Unión: $[-3,2]$.
\end{enumerate}
\resultado \textbf{$[-3,2]$}, que es exactamente el complementario de la solución estricta $(-\infty,-3)\cup(2,\infty)$ de la desigualdad $>5$; los puntos $-3$ y $2$ (suma de distancias igual a~5) pertenecen a ambos conjuntos «$\le5$» y «$\ge5$».
\comprobacion $x=2$: $1+4=5$; $x=-3$: $4+1=5$; $x=0$: $1+2=3$.
\end{solucion}

\begin{solucion}{eje:2-amgm-aplic}
\planteamiento aplicar el \cref{teo:amgm} a números cuyo producto (o suma) es constante.
\pasos
\begin{enumerate}[label=(\alph*)]
  \item $x+\frac4x\ge2\sqrt{x\cdot\frac4x}=4$, con igualdad si $x=\frac4x$, es decir, $x=2$.
  \item Lados $a,b>0$ con $2a+2b=20$, es decir, $a+b=10$. Área $ab\le\bigl(\frac{a+b}2\bigr)^2=25$, con igualdad si $a=b=5$.
  \item $\frac{x+y+z}3\ge\sqrt[3]{xyz}=\sqrt[3]8=2$, luego $x+y+z\ge6$, con igualdad si $x=y=z=2$.
\end{enumerate}
\resultado \textbf{(a) mínimo $4$ en $x=2$; (b) el cuadrado de lado~5, área~25; (c) mínimo $6$ en $x=y=z=2$.}
\comprobacion (a): $x=1$ da $5$ y $x=4$ da $5$, ambos mayores que $4$.
\end{solucion}

\begin{solucion}{eje:2-bernoulli-num}
\planteamiento aplicar el \cref{teo:bernoulli} con $x=0{,}01$, $n=100$ y con $x=-0{,}01$, $n=50$.
\pasos
\begin{enumerate}[label=\arabic*.]
  \item $1{,}01^{100}=(1+0{,}01)^{100}\ge1+100\cdot0{,}01=2$.
  \item $0{,}99^{50}=(1-0{,}01)^{50}\ge1-50\cdot0{,}01=0{,}5$; se puede aplicar porque $-0{,}01\ge-1$.
  \item Valores reales: $1{,}01^{100}\approx2{,}705$ y $0{,}99^{50}\approx0{,}6050$.
\end{enumerate}
\resultado \textbf{$1{,}01^{100}\ge2$ y $0{,}99^{50}\ge0{,}5$}; las cotas son correctas pero no ajustadas.
\comprobacion la cota de Bernoulli es la recta tangente en $x=0$; se aleja del valor real cuando $nx$ no es pequeño.
\end{solucion}

\begin{solucion}{eje:2-cs-num}
\planteamiento (a) evaluar ambos miembros; (b) aplicar el \cref{teo:cs} a $(2,3)$ y $(x,y)$.
\pasos
\begin{enumerate}[label=\arabic*.]
  \item (a) $\sum a_ib_i=4-2+6=8$; $\sum a_i^2=1+4+9=14$; $\sum b_i^2=16+1+4=21$. Se cumple $64\le294$.
  \item (b) $(2x+3y)^2\le(2^2+3^2)(x^2+y^2)=13\cdot13=169$, luego $2x+3y\le13$.
  \item La igualdad en Cauchy--Schwarz se da cuando $(x,y)$ es proporcional a $(2,3)$: $(x,y)=t(2,3)$ con $13t^2=13$; para el máximo, $t=1$.
\end{enumerate}
\resultado \textbf{(a) $64\le294$; (b) máximo $13$, alcanzado en $(2,3)$.}
\comprobacion $2^2+3^2=13$ y $2\cdot2+3\cdot3=13$.
\end{solucion}

\begin{solucion}{eje:2-triangular-n}
\datos $x_1,\dots,x_n\in\R$.
\planteamiento inducción sobre $n$ usando el \cref{teo:triangular} en el paso.
\pasos
\begin{enumerate}[label=\arabic*.]
  \item Caso $n=1$: $\abs{x_1}\le\abs{x_1}$.
  \item Paso: sea $S=x_1+\dots+x_n$ y supóngase $\abs S\le\sum_{i=1}^n\abs{x_i}$.
  \item Por el \cref{teo:triangular}, $\abs{S+x_{n+1}}\le\abs S+\abs{x_{n+1}}$.
  \item Por la hipótesis de inducción, $\abs S+\abs{x_{n+1}}\le\sum_{i=1}^{n+1}\abs{x_i}$.
\end{enumerate}
\resultado \textbf{$\abs{\sum_{i=1}^nx_i}\le\sum_{i=1}^n\abs{x_i}$ para todo $n$.}
\comprobacion $\abs{3-5+1}=1\le9=3+5+1$.
\end{solucion}

\begin{solucion}{eje:2-maxmin}
\planteamiento separar los casos $x\ge y$ y $x<y$.
\pasos
\begin{enumerate}[label=\arabic*.]
  \item Si $x\ge y$: $\abs{x-y}=x-y$; $\frac{x+y+(x-y)}2=x=\max(x,y)$ y $\frac{x+y-(x-y)}2=y=\min(x,y)$.
  \item Si $x<y$: $\abs{x-y}=y-x$; $\frac{x+y+(y-x)}2=y=\max(x,y)$ y $\frac{x+y-(y-x)}2=x=\min(x,y)$.
\end{enumerate}
\resultado \textbf{Las dos fórmulas son ciertas para todos $x,y$.}
\comprobacion $x=2$, $y=7$: $\frac{9+5}2=7$ y $\frac{9-5}2=2$. Además $\max+\min=x+y$, como debe ser.
\end{solucion}

\begin{solucion}{eje:2-cadena-medias}
\datos $0<a<b$.
\planteamiento probar las tres desigualdades por separado.
\pasos
\begin{enumerate}[label=\arabic*.]
  \item $a<\sqrt{ab}$: ambos son positivos y $a^2<ab$ (multiplicando $a<b$ por $a>0$); por el \cref{eje:1-potencia-n}, $a<\sqrt{ab}$.
  \item $\sqrt{ab}<\frac{a+b}2$: por la \cref{prop:amgm2}, con desigualdad estricta porque $a\neq b$.
  \item $\frac{a+b}2<b$: es el \cref{eje:1-epsilon}(a).
\end{enumerate}
\resultado \textbf{$a<\sqrt{ab}<\frac{a+b}2<b$.}
\comprobacion $a=1$, $b=4$: $1<2<2{,}5<4$.
\end{solucion}

\begin{solucion}{eje:2-media-armonica}
\datos $a,b>0$; $G=\sqrt{ab}$, $A=\frac{a+b}2$.
\planteamiento reducir cada desigualdad a la \cref{prop:amgm2} o a un cuadrado no negativo.
\pasos
\begin{enumerate}[label=\arabic*.]
  \item Armónica--geométrica: $\frac{2ab}{a+b}=\frac{G^2}{A}\le\frac{G^2}{G}=G$, porque $A\ge G>0$ (\cref{eje:1-negativo}(c)). Igualdad si y solo si $A=G$, es decir, $a=b$.
  \item Geométrica--aritmética: es la \cref{prop:amgm2}.
  \item Aritmética--cuadrática: ambos miembros son positivos, así que basta comparar cuadrados: $\frac{a^2+b^2}2-\frac{(a+b)^2}4=\frac{2a^2+2b^2-a^2-2ab-b^2}4=\frac{(a-b)^2}4\ge0$, con igualdad si y solo si $a=b$.
\end{enumerate}
\resultado \textbf{$\frac{2ab}{a+b}\le\sqrt{ab}\le\frac{a+b}2\le\sqrt{\frac{a^2+b^2}2}$}, con igualdades si y solo si $a=b$.
\comprobacion $a=1$, $b=4$: $1{,}6\le2\le2{,}5\le2{,}915$.
\end{solucion}

\begin{solucion}{eje:2-bernoulli-fuerte}
\datos $x\ge0$, $n\in\N$.
\planteamiento inducción como en el \cref{teo:bernoulli}, descartando solo términos no negativos.
\pasos
\begin{enumerate}[label=\arabic*.]
  \item $n=1$: $1+x\ge1+x+0$.
  \item Paso: si $(1+x)^n\ge1+nx+\frac{n(n-1)}2x^2$, multiplicando por $1+x>0$:
  \[
  (1+x)^{n+1}\ge1+(n+1)x+\Bigl(n+\frac{n(n-1)}2\Bigr)x^2+\frac{n(n-1)}2x^3.
  \]
  \item $n+\frac{n(n-1)}2=\frac{n(n+1)}2=\frac{(n+1)n}2$ y el término en $x^3$ es $\ge0$; se obtiene el enunciado para $n+1$.
  \item Con $x=1$: $2^n\ge1+n+\frac{n(n-1)}2>\frac{n(n-1)}2$.
\end{enumerate}
\resultado \textbf{$(1+x)^n\ge1+nx+\frac{n(n-1)}2x^2$ para $x\ge0$, y $2^n>\frac{n(n-1)}2$.}
\comprobacion $n=3$, $x=1$: $8\ge1+3+3=7$.
\end{solucion}

\begin{solucion}{eje:2-suma-inversos}
\datos $x_1,\dots,x_n>0$.
\planteamiento aplicar el \cref{teo:cs} a $a_i=\sqrt{x_i}$, $b_i=1/\sqrt{x_i}$.
\pasos
\begin{enumerate}[label=\arabic*.]
  \item $\sum_ia_ib_i=\sum_i1=n$; $\sum_ia_i^2=\sum_ix_i$; $\sum_ib_i^2=\sum_i1/x_i$.
  \item Cauchy--Schwarz: $n^2\le\bigl(\sum_ix_i\bigr)\bigl(\sum_i1/x_i\bigr)$.
  \item Igualdad: si todos los $x_i$ valen $c$, el producto es $nc\cdot\frac nc=n^2$. Recíprocamente, por la desigualdad de las medias, $\sum x_i\ge n\sqrt[n]{\prod x_i}$ y $\sum\frac1{x_i}\ge n\big/\sqrt[n]{\prod x_i}$; si el producto vale $n^2$, ambas son igualdades y, por el \cref{teo:amgm}, todos los $x_i$ son iguales.
\end{enumerate}
\resultado \textbf{$\bigl(\sum x_i\bigr)\bigl(\sum 1/x_i\bigr)\ge n^2$, con igualdad si y solo si todos los $x_i$ coinciden.}
\comprobacion $x=(1,2)$: $3\cdot1{,}5=4{,}5\ge4$.
\end{solucion}
% ---------------------------------------------------------------------
\section{Tema 3}

\begin{solucion}{eje:3-basico}
\datos $P(x)=2x^3-x+5$; raíces pedidas $3$ y $-\frac12$.
\planteamiento evaluar; usar $\deg(PQ)=\deg P+\deg Q$; construir el polinomio a partir de sus factores (\cref{prop:vieta}).
\pasos
\begin{enumerate}[label=\arabic*.]
  \item (a) $P(-1)=2(-1)-(-1)+5=-2+1+5=4$.
  \item (b) $\deg(P^2(x^2+1))=3+3+2=8$.
  \item (c) $(x-3)\bigl(x+\frac12\bigr)=0$; multiplicando por~2 para tener coeficientes enteros: $(x-3)(2x+1)=2x^2-5x-3=0$.
\end{enumerate}
\resultado \textbf{(a) $4$; (b) $8$; (c) $2x^2-5x-3=0$.}
\comprobacion (c) por Vieta: suma $\frac52=3-\frac12$, producto $-\frac32=3\cdot(-\frac12)$.
\end{solucion}

\begin{solucion}{eje:3-discriminante}
\planteamiento calcular $\Delta=b^2-4ac$ (\cref{teo:cuadratica}).
\pasos
\begin{enumerate}[label=(\alph*)]
  \item $\Delta=16-20=-4<0$: ninguna solución real.
  \item $\Delta=144-144=0$: una solución doble, $x=\frac{12}8=\frac32$.
  \item $\Delta=1+24=25>0$: dos soluciones, $x=\frac{-1\pm5}{4}$, es decir, $1$ y $-\frac32$.
\end{enumerate}
\resultado \textbf{(a) ninguna; (b) una (doble); (c) dos.}
\comprobacion (b) $4x^2-12x+9=(2x-3)^2$.
\end{solucion}

\begin{solucion}{eje:3-division}
\datos $P=2x^4-3x^3+0x^2+x-5$, $D=x^2-x+1$.
\planteamiento algoritmo de la demostración del \cref{teo:division}: eliminar en cada paso el término de mayor grado.
\pasos
\begin{enumerate}[label=\arabic*.]
  \item $2x^4/x^2=2x^2$; $P-2x^2D=-x^3-2x^2+x-5$.
  \item $-x^3/x^2=-x$; restando $-xD=-x^3+x^2-x$: queda $-3x^2+2x-5$.
  \item $-3x^2/x^2=-3$; restando $-3D=-3x^2+3x-3$: queda $-x-2$, de grado $1<2$.
\end{enumerate}
\resultado \textbf{$Q=2x^2-x-3$, $R=-x-2$.}
\comprobacion $(x^2-x+1)(2x^2-x-3)=2x^4-3x^3+0x^2+2x-3$; sumando $R$: $2x^4-3x^3+x-5=P$.
\end{solucion}

\begin{solucion}{eje:3-ruffini}
\planteamiento probar divisores del término independiente (\cref{teo:raices-racionales}) y dividir por Ruffini.
\pasos
\begin{enumerate}[label=\arabic*.]
  \item (a) $P(1)=1-6+11-6=0$. Ruffini con $1$: coeficientes $1,-5,6$, es decir, $x^2-5x+6=(x-2)(x-3)$.
  \item (b) Con $t=x^2$: $t^2-5t+4=(t-1)(t-4)$, luego $(x^2-1)(x^2-4)=(x-1)(x+1)(x-2)(x+2)$. Por Ruffini con $x=1$ (coeficientes $1,0,-5,0,4$): cociente $x^3+x^2-4x-4$, resto~0; y así sucesivamente.
\end{enumerate}
\resultado \textbf{(a) $(x-1)(x-2)(x-3)$; (b) $(x-1)(x+1)(x-2)(x+2)$.}
\comprobacion (a) en $x=0$: $(-1)(-2)(-3)=-6$, igual al término independiente.
\end{solucion}

\begin{solucion}{eje:3-resto}
\planteamiento teorema del resto (\cref{teo:resto}): el resto entre $x-a$ es $P(a)$.
\pasos
\begin{enumerate}[label=\arabic*.]
  \item (a) $P(2)=8+4k-8+4=4k+4=0$, luego $k=-1$.
  \item (b) $x+1=x-(-1)$: el resto de $x^{100}-1$ es $(-1)^{100}-1=0$; el de $x^{101}+3$ es $(-1)^{101}+3=2$.
\end{enumerate}
\resultado \textbf{(a) $k=-1$; (b) restos $0$ y $2$.}
\comprobacion (a) $x^3-x^2-4x+4=(x-1)(x^2-4)=(x-1)(x-2)(x+2)$, divisible por $x-2$.
\end{solucion}

\begin{solucion}{eje:3-racionales}
\datos $P(x)=2x^3-3x^2-3x+2$.
\planteamiento por el \cref{teo:raices-racionales}, las raíces racionales son de la forma $p/q$ con $p\mid2$, $q\mid2$: candidatas $\pm1,\pm2,\pm\frac12$.
\pasos
\begin{enumerate}[label=\arabic*.]
  \item $P(1)=-2$, $P(-1)=0$. Ruffini con $-1$: coeficientes $2,-5,2$.
  \item $2x^2-5x+2=0$: $x=\frac{5\pm3}4$, es decir, $2$ y $\frac12$.
\end{enumerate}
\resultado \textbf{Raíces $-1$, $\frac12$ y $2$; $P(x)=(x+1)(2x-1)(x-2)$.}
\comprobacion Vieta: suma $-1+\frac12+2=\frac32=-\frac{-3}2$; producto $-1=(-1)^3\cdot\frac22$.
\end{solucion}

\begin{solucion}{eje:3-vieta}
\datos $x_1+x_2=5$, $x_1x_2=3$ (\cref{prop:vieta}).
\planteamiento expresar cada cantidad mediante la suma $s$ y el producto $p$.
\pasos
\begin{enumerate}[label=\arabic*.]
  \item $x_1^2+x_2^2=s^2-2p=25-6=19$.
  \item $\frac1{x_1}+\frac1{x_2}=\frac{s}{p}=\frac53$.
  \item $x_1^3+x_2^3=s^3-3ps=125-45=80$.
\end{enumerate}
\resultado \textbf{$19$, $\frac53$ y $80$.}
\comprobacion las raíces son $\frac{5\pm\sqrt{13}}2\approx4{,}303$ y $0{,}6972$; $4{,}303^2+0{,}6972^2\approx18{,}52+0{,}486\approx19{,}0$.
\end{solucion}

\begin{solucion}{eje:3-ecuaciones}
\planteamiento (a) cambio $t=x^2$; (b) elevar al cuadrado y comprobar; (c) multiplicar por $(x-1)(x+1)$ con $x\neq\pm1$.
\pasos
\begin{enumerate}[label=\arabic*.]
  \item (a) $t^2-13t+36=0$, $t=4$ o $t=9$; $x=\pm2$ o $x=\pm3$.
  \item (b) $x+5=(x-1)^2=x^2-2x+1$, luego $x^2-3x-4=0$ y $x=4$ o $x=-1$. Comprobación: $x=4$ da $\sqrt9=3=4-1$, válida; $x=-1$ da $\sqrt4=2\neq-2$, extraña.
  \item (c) $(x+1)+2(x-1)=x^2-1$, es decir, $x^2-3x=0$, $x=0$ o $x=3$; ninguna anula los denominadores.
\end{enumerate}
\resultado \textbf{(a) $\pm2,\pm3$; (b) $x=4$; (c) $x=0$ o $x=3$.}
\comprobacion (c) $x=3$: $\frac12+\frac24=1$; $x=0$: $-1+2=1$.
\end{solucion}

\begin{solucion}{eje:3-inecuaciones}
\planteamiento factorizar y hacer la tabla de signos (\cref{ej:tabla-signos}).
\pasos
\begin{enumerate}[label=\arabic*.]
  \item (a) $x^2-x-6=(x-3)(x+2)$; es negativo entre las raíces: $(-2,3)$.
  \item (b) Puntos críticos $-2$, $1$, $3$. Signos de $\frac{(x-1)(x+2)}{x-3}$: en $(-\infty,-2)$, $\frac{(-)(-)}{(-)}<0$; en $(-2,1)$, $\frac{(-)(+)}{(-)}>0$; en $(1,3)$, $\frac{(+)(+)}{(-)}<0$; en $(3,\infty)$, $>0$.
  \item Se incluyen los ceros del numerador ($-2$ y $1$) y se excluye $3$.
\end{enumerate}
\resultado \textbf{(a) $(-2,3)$; (b) $[-2,1]\cup(3,\infty)$.}
\comprobacion (b) $x=0$: $\frac{(-1)(2)}{-3}=\frac23\ge0$; $x=2$: $\frac{1\cdot4}{-1}<0$.
\end{solucion}

\begin{solucion}{eje:3-sistemas}
\planteamiento (a) regla de Cramer (\cref{teo:cramer}) y estudio aparte de $\delta=0$; (b) reducción eliminando primero $z$.
\pasos
\begin{enumerate}[label=\arabic*.]
  \item (a) $\delta=1\cdot k-2\cdot3=k-6$. Si $k\neq6$: $x=\frac{4k-2\cdot6}{k-6}=\frac{4k-12}{k-6}$, $y=\frac{1\cdot6-4\cdot3}{k-6}=\frac{-6}{k-6}$.
  \item Si $k=6$: la segunda ecuación es $3x+6y=6$, es decir, $x+2y=2$, incompatible con $x+2y=4$.
  \item (b) Sumando la primera y la tercera: $2x+3y=8$. Sumando la segunda y la tercera: $3x+y=5$.
  \item De la última, $y=5-3x$; sustituyendo, $2x+15-9x=8$, $x=1$, $y=2$, y de la primera $z=3$.
\end{enumerate}
\resultado \textbf{(a) si $k\neq6$, $(x,y)=\bigl(\frac{4k-12}{k-6},\frac{-6}{k-6}\bigr)$; si $k=6$, incompatible. (b) $(1,2,3)$.}
\comprobacion (a) $k=0$: $(2,1)$, y $2+2=4$, $6+0=6$. (b) coincide con el \cref{ej:gauss3}.
\end{solucion}

\begin{solucion}{eje:3-identidad}
\datos $\deg P,\deg Q\le n$; $P(c_i)=Q(c_i)$ para $n+1$ valores distintos $c_0,\dots,c_n$.
\planteamiento aplicar el \cref{cor:raices} a $R=P-Q$.
\pasos
\begin{enumerate}[label=\arabic*.]
  \item $R=P-Q$ es nulo o tiene grado a lo sumo $n$.
  \item $R(c_i)=0$ para $n+1$ valores distintos.
  \item Si $R$ no fuera nulo, tendría como mucho $n$ raíces (\cref{cor:raices}): contradicción. Luego $R=0$ y $P=Q$.
\end{enumerate}
\resultado \textbf{$P=Q$ coeficiente a coeficiente.}
\comprobacion el número $n+1$ no puede rebajarse: $P=x^2$ y $Q=x$ ($n=2$) coinciden en los dos puntos $0$ y $1$, y son distintos.
\end{solucion}

\begin{solucion}{eje:3-diferencia-potencias}
\datos $n\in\N$.
\planteamiento desarrollar el producto: la suma es telescópica.
\pasos
\begin{enumerate}[label=\arabic*.]
  \item $(x-a)\sum_{k=0}^{n-1}a^kx^{n-1-k}=\sum_{k=0}^{n-1}a^kx^{n-k}-\sum_{k=0}^{n-1}a^{k+1}x^{n-1-k}$.
  \item Con $j=k+1$ en la segunda suma, esta es $\sum_{j=1}^{n}a^jx^{n-j}$. Al restar se cancelan los términos con $1\le j\le n-1$ y queda $x^n-a^n$ (\cref{prop:telescopica}).
  \item Por tanto $x-a$ divide a $x^n-a^n$.
  \item Si $n$ es impar, sustituyendo $a$ por $-a$: $x^n-(-a)^n=x^n+a^n$ es divisible por $x-(-a)=x+a$.
\end{enumerate}
\resultado \textbf{La identidad es cierta; $x-a\mid x^n-a^n$ siempre y $x+a\mid x^n+a^n$ para $n$ impar.}
\comprobacion $n=3$: $x^3+8=(x+2)(x^2-2x+4)$.
\end{solucion}

\begin{solucion}{eje:3-impares}
\datos $b,c$ impares.
\planteamiento por el \cref{teo:raices-racionales} con $a_n=1$, toda raíz racional es entera; se estudia la paridad.
\pasos
\begin{enumerate}[label=\arabic*.]
  \item Si $p/q$ es raíz racional irreducible, $q\mid1$, luego la raíz es un entero $r$.
  \item Si $r$ es par: $r^2$ y $br$ son pares y $c$ impar, así que $r^2+br+c$ es impar.
  \item Si $r$ es impar: $r^2$, $br$ y $c$ son impares y su suma es impar.
  \item En ambos casos $r^2+br+c\neq0$, contradicción.
\end{enumerate}
\resultado \textbf{$x^2+bx+c=0$ no tiene raíces racionales.}
\comprobacion $x^2+3x+1=0$ tiene raíces $\frac{-3\pm\sqrt5}2$, irracionales.
\end{solucion}

\begin{solucion}{eje:3-det-cero}
\datos $ad-bc=0$; algún coeficiente no nulo; por simetría se supone $a\neq0$.
\planteamiento mostrar que la segunda ecuación, en su parte izquierda, es múltiplo de la primera.
\pasos
\begin{enumerate}[label=\arabic*.]
  \item Sea $\lambda=c/a$. Entonces $c=\lambda a$ y $d=\frac{bc}a=\lambda b$.
  \item La segunda ecuación es $\lambda(ax+by)=f$. Si $(x,y)$ resuelve la primera, debe cumplirse $\lambda e=f$.
  \item Si $f\neq\lambda e$, no hay solución.
  \item Si $f=\lambda e$, la segunda ecuación es consecuencia de la primera, cuyas soluciones son $x=\frac{e-by}a$ con $y$ libre: infinitas.
\end{enumerate}
\resultado \textbf{Con $\delta=0$ el sistema es incompatible o compatible indeterminado.}
\comprobacion el sistema del \cref{eje:3-sistemas} con $k=6$ tiene $\lambda=3$, $\lambda e=12\neq6=f$: incompatible.
\end{solucion}

\begin{solucion}{eje:3-cuadratica-positiva}
\datos $a>0$, $f(x)=ax^2+bx+c$.
\planteamiento completar el cuadrado como en el \cref{teo:cuadratica}.
\pasos
\begin{enumerate}[label=\arabic*.]
  \item $f(x)=a\bigl(x+\frac b{2a}\bigr)^2+c-\frac{b^2}{4a}$.
  \item Como $a>0$ y un cuadrado es $\ge0$, $f(x)\ge c-\frac{b^2}{4a}$, con igualdad si y solo si $x=-\frac b{2a}$.
  \item $f\ge0$ en todo $\R$ si y solo si su mínimo es $\ge0$: $c-\frac{b^2}{4a}\ge0$; multiplicando por $4a>0$, $4ac-b^2\ge0$.
\end{enumerate}
\resultado \textbf{Mínimo $c-\frac{b^2}{4a}$ en $x=-\frac b{2a}$; $f\ge0$ si y solo si $b^2-4ac\le0$.}
\comprobacion $f=x^2-2x+5$: mínimo en $x=1$ con valor $4=5-\frac44$, y $\Delta=-16\le0$.
\end{solucion}
% ---------------------------------------------------------------------
\section{Tema 4}

\begin{solucion}{eje:4-notacion}
\planteamiento encontrar el término general y los límites del índice.
\pasos
\begin{enumerate}[label=(\alph*)]
  \item Impares $2k-1$; $2k-1=99$ da $k=50$: $\sum_{k=1}^{50}(2k-1)$.
  \item Potencias $2^k$; $2^{10}=1024$: $\sum_{k=1}^{10}2^k$.
  \item Términos $(-1)^k/k$ para $k=2,\dots,100$: $\sum_{k=2}^{100}\frac{(-1)^k}{k}$.
\end{enumerate}
\resultado \textbf{(a) $\sum_{k=1}^{50}(2k-1)$; (b) $\sum_{k=1}^{10}2^k$; (c) $\sum_{k=2}^{100}(-1)^k/k$.}
\comprobacion (a) por el \cref{teo:aritmetica}, $\frac{50(1+99)}2=2500=50^2$, la conocida suma de los primeros impares.
\end{solucion}

\begin{solucion}{eje:4-calcular}
\planteamiento evaluar término a término y aplicar los convenios de la suma vacía y de la suma de una constante.
\pasos
\begin{enumerate}[label=(\alph*)]
  \item $1+4+9+16=30$.
  \item $\frac12\cdot\frac23\cdot\frac34\cdot\frac45\cdot\frac56=\frac16$: cada numerador se cancela con el denominador anterior.
  \item El límite superior es menor que el inferior: suma vacía, que vale~$0$.
  \item $\sum_{k=1}^{10}7=10\cdot7=70$ (\cref{prop:sumas}(iii)).
\end{enumerate}
\resultado \textbf{(a) $30$; (b) $1/6$; (c) $0$; (d) $70$.}
\comprobacion (a) con el \cref{teo:potencias}: $\frac{4\cdot5\cdot9}6=30$.
\end{solucion}

\begin{solucion}{eje:4-errores}
\planteamiento comparar con la \cref{prop:sumas} y la \cref{prop:producto-sumas}.
\pasos
\begin{enumerate}[label=(\alph*)]
  \item $\sum_{k=1}^n3=3n$: se suman $n$ treses.
  \item $(a_1+a_2)(b_1+b_2)=a_1b_1+a_2b_2+a_1b_2+a_2b_1$: faltan los términos cruzados. Contraejemplo: $a=b=(1,1)$ da $2$ frente a $4$.
  \item $(a_1+a_2)^2=a_1^2+a_2^2+2a_1a_2$: falta el doble producto.
\end{enumerate}
\resultado \textbf{(a) vale $3n$; (b) y (c) olvidan los términos cruzados.}
\comprobacion (c) con $a=(1,2)$: $9\neq5$, y $5+2\cdot2=9$.
\end{solucion}

\begin{solucion}{eje:4-aritmeticas}
\planteamiento linealidad y \cref{teo:aritmetica}.
\pasos
\begin{enumerate}[label=\arabic*.]
  \item (a) $\sum(3k-2)=3\cdot\frac{100\cdot101}2-2\cdot100=15150-200=14950$.
  \item (b) El primer múltiplo de~7 mayor que 100 es $105=7\cdot15$; el último menor que 1000 es $994=7\cdot142$.
  \item Hay $142-15+1=128$ términos. Suma: $\frac{128(105+994)}2=64\cdot1099=70336$.
\end{enumerate}
\resultado \textbf{(a) $14950$; (b) $70336$.}
\comprobacion (b) $7\sum_{k=15}^{142}k=7\bigl(\frac{142\cdot143}2-\frac{14\cdot15}2\bigr)=7(10153-105)=70336$.
\end{solucion}

\begin{solucion}{eje:4-geometricas}
\planteamiento \cref{teo:geometrica}.
\pasos
\begin{enumerate}[label=(\alph*)]
  \item $3\sum_{k=0}^92^k=3\cdot\frac{2^{10}-1}{2-1}=3\cdot1023=3069$.
  \item $\sum_{k=1}^n(\frac12)^k=\sum_{k=0}^n(\frac12)^k-1=\frac{1-2^{-(n+1)}}{1/2}-1=2-2^{-n}-1=1-2^{-n}$.
  \item Razón $r=-x\neq1$: $\sum_{k=0}^n(-x)^k=\frac{1-(-x)^{n+1}}{1+x}$.
\end{enumerate}
\resultado \textbf{(a) $3069$; (b) $1-2^{-n}$; (c) $\dfrac{1-(-x)^{n+1}}{1+x}$.}
\comprobacion (b) $n=2$: $\frac12+\frac14=\frac34=1-\frac14$. (c) $n=1$: $1-x=\frac{1-x^2}{1+x}$.
\end{solucion}

\begin{solucion}{eje:4-telescopicas}
\planteamiento escribir cada término como diferencia y aplicar la \cref{prop:telescopica}.
\pasos
\begin{enumerate}[label=\arabic*.]
  \item (a) $\frac1{k(k+2)}=\frac12\bigl(\frac1k-\frac1{k+2}\bigr)$. Al sumar sobreviven $\frac11+\frac12$ con signo $+$ y $\frac1{n+1}+\frac1{n+2}$ con signo $-$: $\frac12\bigl(\frac32-\frac1{n+1}-\frac1{n+2}\bigr)$.
  \item Para $n=10$: $\frac12\bigl(\frac32-\frac1{11}-\frac1{12}\bigr)=\frac12\cdot\frac{198-12-11}{132}=\frac{175}{264}$.
  \item (b) $\ln(1+\frac1k)=\ln(k+1)-\ln k$; suma $\ln(n+1)-\ln1=\ln(n+1)$.
  \item (c) $\frac1{\sqrt k+\sqrt{k+1}}=\frac{\sqrt{k+1}-\sqrt k}{(k+1)-k}=\sqrt{k+1}-\sqrt k$; suma $\sqrt{100}-\sqrt1=9$.
\end{enumerate}
\resultado \textbf{(a) $\frac34-\frac{2n+3}{2(n+1)(n+2)}$, que vale $\frac{175}{264}$ para $n=10$; (b) $\ln(n+1)$; (c) $9$.}
\comprobacion (a) $n=1$: $\frac13$ y $\frac34-\frac5{12}=\frac4{12}=\frac13$.
\end{solucion}

\begin{solucion}{eje:4-indice}
\planteamiento (a) cambio $j=k-2$ (\cref{prop:sumas}(iv)); (b) $k(k+1)=k^2+k$ y \cref{teo:potencias}.
\pasos
\begin{enumerate}[label=\arabic*.]
  \item (a) Si $k=3$, $j=1$; si $k=n+2$, $j=n$: $\sum_{j=1}^nj^2=\frac{n(n+1)(2n+1)}6$.
  \item (b) $\sum_{k=1}^n(k^2+k)=\frac{n(n+1)(2n+1)}6+\frac{n(n+1)}2=\frac{n(n+1)(2n+1+3)}6=\frac{n(n+1)(n+2)}3$.
  \item Para $n=10$: $\frac{10\cdot11\cdot12}3=440$.
\end{enumerate}
\resultado \textbf{(a) $\sum_{j=1}^nj^2=\frac{n(n+1)(2n+1)}6$; (b) $\frac{n(n+1)(n+2)}3$, que vale $440$ para $n=10$.}
\comprobacion (b) $n=10$ directamente: $385+55=440$.
\end{solucion}

\begin{solucion}{eje:4-dobles}
\planteamiento sumar primero el índice interior; en (c), reescribir el dominio triangular.
\pasos
\begin{enumerate}[label=\arabic*.]
  \item (a) $\sum_{j=1}^4(i+2j)=4i+2\cdot10=4i+20$; $\sum_{i=1}^3(4i+20)=24+60=84$.
  \item (b) $\sum_{j=1}^ij=\frac{i(i+1)}2$; $\sum_{i=1}^n\frac{i(i+1)}2=\frac{n(n+1)(n+2)}6$ por el \cref{eje:4-indice}.
  \item (c) El dominio es $1\le i\le j\le n$: $\sum_{j=1}^n\sum_{i=1}^j\frac1j=\sum_{j=1}^n\frac jj=n$.
\end{enumerate}
\resultado \textbf{(a) $84$; (b) $\frac{n(n+1)(n+2)}6$; (c) $n$.}
\comprobacion (c) $n=2$: $(\frac11+\frac12)+\frac12=2$.
\end{solucion}

\begin{solucion}{eje:4-binomio}
\planteamiento \cref{teo:binomio} con los valores adecuados de $a$ y $b$.
\pasos
\begin{enumerate}[label=\arabic*.]
  \item (a) Con $a=2x$, $b=-1$ y $\binom5k=1,5,10,10,5,1$: $32x^5-80x^4+80x^3-40x^2+10x-1$.
  \item (b) Término general $\binom8kx^{8-k}(2/x)^k=\binom8k2^kx^{8-2k}$; $8-2k=4$ da $k=2$: $\binom82\cdot4=28\cdot4=112$.
  \item (c) $\sum_k\binom6k1^{6-k}2^k=(1+2)^6=729$.
\end{enumerate}
\resultado \textbf{(a) $32x^5-80x^4+80x^3-40x^2+10x-1$; (b) $112$; (c) $729$.}
\comprobacion (a) en $x=1$: $32-80+80-40+10-1=1=(2-1)^5$.
\end{solucion}

\begin{solucion}{eje:4-productos}
\planteamiento factorizar y cancelar (producto telescópico); exponentes con el \cref{teo:aritmetica}.
\pasos
\begin{enumerate}[label=\arabic*.]
  \item (a) $1-\frac1{k^2}=\frac{k-1}k\cdot\frac{k+1}k$. Entonces $\prod_{k=2}^n\frac{k-1}k=\frac1n$ y $\prod_{k=2}^n\frac{k+1}k=\frac{n+1}2$; el producto es $\frac{n+1}{2n}$, que para $n=10$ vale $\frac{11}{20}$.
  \item (b) $\prod2^k=2^{\sum k}=2^{n(n+1)/2}$.
  \item (c) $\ln\prod\frac{k+1}k=\sum(\ln(k+1)-\ln k)=\ln(n+1)$.
\end{enumerate}
\resultado \textbf{(a) $\frac{n+1}{2n}$ ($\frac{11}{20}$ para $n=10$); (b) $2^{n(n+1)/2}$; (c) $\ln(n+1)$.}
\comprobacion (a) $n=2$: $1-\frac14=\frac34=\frac{2+1}{2\cdot2}$.
\end{solucion}

\begin{solucion}{eje:4-cubos}
\datos $S_3(n)=\sum_{k=1}^nk^3$.
\planteamiento inducción sobre $n$ (\cref{teo:induccion}).
\pasos
\begin{enumerate}[label=\arabic*.]
  \item $n=1$: $1=\bigl(\frac{1\cdot2}2\bigr)^2$.
  \item Paso: si $S_3(n)=\frac{n^2(n+1)^2}4$, entonces $S_3(n+1)=\frac{n^2(n+1)^2}4+(n+1)^3=\frac{(n+1)^2(n^2+4n+4)}4$.
  \item Como $n^2+4n+4=(n+2)^2$, $S_3(n+1)=\bigl(\frac{(n+1)(n+2)}2\bigr)^2$, que es el enunciado para $n+1$.
\end{enumerate}
\resultado \textbf{$\sum_{k=1}^nk^3=\bigl(\frac{n(n+1)}2\bigr)^2$ para todo $n$.}
\comprobacion $n=4$: $1+8+27+64=100=10^2$.
\end{solucion}

\begin{solucion}{eje:4-perturbacion}
\datos $S=\sum_{k=0}^nkr^k$, $r\neq1$.
\planteamiento método de perturbación: calcular $S_{n+1}$ de dos formas~\cite[§2.3]{gkp}.
\pasos
\begin{enumerate}[label=\arabic*.]
  \item Separando el último término: $S+(n+1)r^{n+1}=\sum_{k=0}^{n+1}kr^k$.
  \item Separando el primero ($k=0$, que vale~0) y cambiando $k=j+1$: $\sum_{k=0}^{n+1}kr^k=\sum_{j=0}^n(j+1)r^{j+1}=rS+r\sum_{j=0}^nr^j$.
  \item Igualando y usando el \cref{teo:geometrica}: $S(1-r)=r\frac{1-r^{n+1}}{1-r}-(n+1)r^{n+1}$.
  \item Dividiendo por $1-r$ y reduciendo a común denominador:
  \[
  S=\frac{r-r^{n+2}-(n+1)r^{n+1}+(n+1)r^{n+2}}{(1-r)^2}=\frac{r-(n+1)r^{n+1}+nr^{n+2}}{(1-r)^2}.
  \]
\end{enumerate}
\resultado \textbf{$\sum_{k=0}^nkr^k=\dfrac{r-(n+1)r^{n+1}+nr^{n+2}}{(1-r)^2}$.}
\comprobacion $n=3$, $r=2$: $0+2+8+24=34$ y $\frac{2-4\cdot16+3\cdot32}{1}=34$.
\end{solucion}

\begin{solucion}{eje:4-kbinom}
\datos $1\le k\le n$.
\planteamiento desarrollar los factoriales; después, linealidad y \cref{cor:binomio}.
\pasos
\begin{enumerate}[label=\arabic*.]
  \item $k\binom nk=\frac{k\,n!}{k!(n-k)!}=\frac{n\,(n-1)!}{(k-1)!\,((n-1)-(k-1))!}=n\binom{n-1}{k-1}$.
  \item El término $k=0$ vale~0, así que $\sum_{k=0}^nk\binom nk=n\sum_{k=1}^n\binom{n-1}{k-1}$.
  \item Con $j=k-1$: $n\sum_{j=0}^{n-1}\binom{n-1}j=n\,2^{n-1}$.
\end{enumerate}
\resultado \textbf{$k\binom nk=n\binom{n-1}{k-1}$ y $\sum_kk\binom nk=n2^{n-1}$.}
\comprobacion $n=3$: $0+3+6+3=12=3\cdot4$.
\end{solucion}

\begin{solucion}{eje:4-basilea}
\datos $n\in\N$.
\planteamiento inducción; el paso se reduce a una desigualdad entre fracciones.
\pasos
\begin{enumerate}[label=\arabic*.]
  \item $n=1$: $1\le2-1$.
  \item Paso: si $\sum_{k=1}^n\frac1{k^2}\le2-\frac1n$, entonces $\sum_{k=1}^{n+1}\frac1{k^2}\le2-\frac1n+\frac1{(n+1)^2}$.
  \item Basta ver que $-\frac1n+\frac1{(n+1)^2}\le-\frac1{n+1}$, es decir, $\frac1{n+1}+\frac1{(n+1)^2}\le\frac1n$.
  \item El miembro izquierdo es $\frac{n+2}{(n+1)^2}$, y $\frac{n+2}{(n+1)^2}\le\frac1n$ equivale a $n^2+2n\le n^2+2n+1$, que es cierto.
\end{enumerate}
\resultado \textbf{$\sum_{k=1}^n\frac1{k^2}\le2-\frac1n$ para todo $n$}; en particular la suma está acotada por~2.
\comprobacion $n=3$: $1+\frac14+\frac19=1{,}3611\le1{,}6667$.
\end{solucion}

\begin{solucion}{eje:4-varianza}
\datos $a_1,\dots,a_n\in\R$; $S_1=\sum a_i$, $S_2=\sum a_i^2$.
\planteamiento desarrollar el cuadrado y aplicar la \cref{prop:producto-sumas}.
\pasos
\begin{enumerate}[label=\arabic*.]
  \item $\sum_i\sum_j(a_i-a_j)^2=\sum_i\sum_ja_i^2-2\sum_i\sum_ja_ia_j+\sum_i\sum_ja_j^2$.
  \item El primer término es $nS_2$ (no depende de $j$), el tercero también es $nS_2$, y el central es $2S_1^2$.
  \item Total: $2nS_2-2S_1^2$.
  \item El miembro izquierdo es suma de cuadrados, luego $\ge0$: $S_1^2\le nS_2$.
  \item Igualdad si y solo si todos los $(a_i-a_j)^2$ son nulos, es decir, todos los $a_i$ son iguales.
\end{enumerate}
\resultado \textbf{$\sum_{i,j}(a_i-a_j)^2=2nS_2-2S_1^2\ge0$ y $\bigl(\sum a_i\bigr)^2\le n\sum a_i^2$.}
\comprobacion $a=(1,2,3)$: los pares ordenados dan $2(1+4+1)=12$, y $2\cdot3\cdot14-2\cdot36=84-72=12$. Dividiendo por $2n^2$ se obtiene que la varianza $\frac{S_2}n-\bigl(\frac{S_1}n\bigr)^2$ es no negativa (módulo~5).
\end{solucion}
% ---------------------------------------------------------------------
\section{Tema 5}

\begin{solucion}{eje:5-calcular}
\planteamiento escribir el argumento como potencia de la base (\cref{def:log}).
\pasos
\begin{enumerate}[label=\arabic*.]
  \item $64=2^6$, luego $\log_264=6$. $\frac1{81}=3^{-4}$, luego $\log_3\frac1{81}=-4$. $0{,}001=10^{-3}$, luego $\log_{10}0{,}001=-3$.
  \item $\ln e^7=7$.
  \item $8=2^3=(4^{1/2})^3=4^{3/2}$, luego $\log_48=\frac32$.
  \item $10^{\log_{10}5}=5$ por definición.
\end{enumerate}
\resultado \textbf{$6$, $-4$, $-3$, $7$, $\frac32$ y $5$.}
\comprobacion $\log_48$ por cambio de base: $\frac{\log_28}{\log_24}=\frac32$.
\end{solucion}

\begin{solucion}{eje:5-vf}
\planteamiento contrastar con el \cref{teo:log}.
\pasos
\begin{enumerate}[label=(\alph*)]
  \item Falsa: con $a=b=1$, $\ln2\neq0=\ln1+\ln1$.
  \item Verdadera: $\log_2(8^x)=x\log_28=3x$.
  \item Verdadera: $\ln\frac1x=\ln1-\ln x=-\ln x$.
  \item Falsa: con $x=e$, $(\ln e)^2=1\neq2$. Lo correcto es $\ln(x^2)=2\ln x$.
\end{enumerate}
\resultado \textbf{(a) falsa; (b) verdadera; (c) verdadera; (d) falsa.}
\comprobacion (d) con $x=1$ ambos miembros valen 0: un caso particular no basta para afirmar una identidad.
\end{solucion}

\begin{solucion}{eje:5-unidades}
\planteamiento $1\ \text{nat}=1/\ln2$ bits; información $-\log_2p$ bits; para numerar $K$ objetos hacen falta $\lceil\log_2K\rceil$ bits (\cref{prop:cifras}).
\pasos
\begin{enumerate}[label=(\alph*)]
  \item $8/\ln2=8/0{,}6931\approx11{,}54$ bits.
  \item $-\log_2(2^{-10})=10$ bits; en nats, $10\ln2\approx6{,}931$.
  \item $2^{15}=32\,768<50\,257\le65\,536=2^{16}$: hacen falta $16$ bits.
\end{enumerate}
\resultado \textbf{(a) $11{,}54$ bits; (b) $10$ bits $=6{,}931$ nats; (c) $16$ bits.}
\comprobacion $\log_250\,257\approx15{,}62$, cuyo entero superior es~16.
\end{solucion}

\begin{solucion}{eje:5-simplificar}
\planteamiento \cref{teo:log}.
\pasos
\begin{enumerate}[label=(\alph*)]
  \item $\log_2\frac{12}3+\log_22^{-3}=\log_24-3=2-3=-1$.
  \item $\frac12\ln16=\ln4$; $\ln\frac e4=1-\ln4$; suma $1$.
  \item $\log_23\cdot\log_38=\log_23\cdot\frac{\log_28}{\log_23}=3$.
  \item $\log_530=\frac{\ln30}{\ln5}\approx\frac{3{,}401}{1{,}609}\approx2{,}113$.
\end{enumerate}
\resultado \textbf{(a) $-1$; (b) $1$; (c) $3$; (d) $\ln30/\ln5\approx2{,}113$.}
\comprobacion (d) $5^2=25<30<125=5^3$, luego el valor está entre 2 y 3.
\end{solucion}

\begin{solucion}{eje:5-exponenciales}
\planteamiento (a) tomar $\ln$; (b) y (c) cambio de variable $t=2^x$ o $t=e^x$, con $t>0$.
\pasos
\begin{enumerate}[label=\arabic*.]
  \item (a) $(x+1)\ln2=x\ln3$, luego $x(\ln3-\ln2)=\ln2$ y $x=\frac{\ln2}{\ln3-\ln2}\approx1{,}710$.
  \item (b) $t^2-3t-4=(t-4)(t+1)=0$; $t=-1$ se descarta porque $t>0$; $2^x=4$, $x=2$.
  \item (c) $t^2-5t+6=(t-2)(t-3)=0$; $x=\ln2$ o $x=\ln3$.
\end{enumerate}
\resultado \textbf{(a) $x=\frac{\ln2}{\ln(3/2)}\approx1{,}710$; (b) $x=2$; (c) $x=\ln2\approx0{,}6931$ o $x=\ln3\approx1{,}099$.}
\comprobacion (a) $2^{2{,}710}\approx6{,}541\approx3^{1{,}710}$.
\end{solucion}

\begin{solucion}{eje:5-logaritmicas}
\planteamiento agrupar logaritmos, igualar argumentos y comprobar el dominio (\cref{ej:ecuacion-log}).
\pasos
\begin{enumerate}[label=\arabic*.]
  \item (a) Dominio $x>2$. $\log_2(x(x-2))=3$, $x^2-2x-8=0$, $x=4$ o $x=-2$; solo $x=4$ está en el dominio.
  \item (b) Dominio $x>1$. $\ln\frac{x+1}{x-1}=\ln3$, $x+1=3x-3$, $x=2$.
\end{enumerate}
\resultado \textbf{(a) $x=4$; (b) $x=2$.}
\comprobacion (a) $\log_24+\log_22=2+1=3$; (b) $\ln3-\ln1=\ln3$.
\end{solucion}

\begin{solucion}{eje:5-inecuaciones}
\planteamiento usar la monotonía (\cref{teo:log}(iv)): con base menor que 1 el sentido se invierte.
\pasos
\begin{enumerate}[label=\arabic*.]
  \item (a) $0{,}5^x=2^{-x}>8=2^3$; como $2^t$ es creciente, $-x>3$, es decir, $x<-3$.
  \item (b) Dominio $2x-1>0$, $x>\frac12$. Como $\log_3$ es creciente, $2x-1\le3^2=9$, $x\le5$.
\end{enumerate}
\resultado \textbf{(a) $(-\infty,-3)$; (b) $\bigl(\frac12,5\bigr]$.}
\comprobacion (a) $x=-4$: $0{,}5^{-4}=16>8$. (b) $x=5$: $\log_39=2$.
\end{solucion}

\begin{solucion}{eje:5-lse}
\datos $z=(2,1,0,-1)$, $M=2$, $n=4$.
\planteamiento cálculo estable: $\LSE(z)=M+\ln\sum_ie^{z_i-M}$ (\cref{prop:lse}).
\pasos
\begin{enumerate}[label=\arabic*.]
  \item $\sum_ie^{z_i-2}=1+e^{-1}+e^{-2}+e^{-3}\approx1+0{,}3679+0{,}1353+0{,}0498=1{,}553$.
  \item $\LSE(z)\approx2+\ln1{,}553\approx2+0{,}4402=2{,}440$.
  \item $\softmax(z)_i=e^{z_i-2}/1{,}553$: $(0{,}6439;\ 0{,}2369;\ 0{,}0871;\ 0{,}0321)$.
  \item Cotas: $2\le2{,}440\le2+\ln4\approx3{,}386$.
  \item (b) $M=1000$: $1000+\ln(1+e^{-1}+e^{-2})\approx1000+\ln1{,}503=1000{,}408$.
\end{enumerate}
\resultado \textbf{(a) $\LSE\approx2{,}440$, softmax $\approx(0{,}6439;0{,}2369;0{,}0871;0{,}0321)$; (b) $\approx1000{,}408$.}
\comprobacion las componentes de la softmax suman $1{,}0000$.
\end{solucion}

\begin{solucion}{eje:5-perplejidad}
\datos log-verosimilitud media $-2{,}3$ nats por token, es decir, $H=2{,}3$ nats.
\planteamiento cambio de base y definición de perplejidad $e^{H}$ ($H$ en nats) o $2^{H}$ ($H$ en bits).
\pasos
\begin{enumerate}[label=\arabic*.]
  \item (a) $2{,}3/\ln2\approx3{,}318$ bits por token.
  \item (b) $e^{2{,}3}\approx9{,}974$; $2^{3{,}318}=e^{3{,}318\ln2}=e^{2{,}3}$, el mismo valor.
  \item (c) Si cada token tiene probabilidad $q$, la log-verosimilitud media es $\ln q=-2{,}3$, luego $q=e^{-2{,}3}\approx0{,}1003$.
\end{enumerate}
\resultado \textbf{(a) $3{,}318$ bits/token; (b) perplejidad $\approx9{,}97$; (c) $q\approx0{,}100$.}
\comprobacion la perplejidad es $1/q$: $1/0{,}1003\approx9{,}97$, el modelo duda como si eligiera entre unos 10 tokens equiprobables.
\end{solucion}

\begin{solucion}{eje:5-ley-potencia}
\datos $L(N)=(N_c/N)^{0{,}076}$.
\planteamiento cociente $L(\lambda N)/L(N)=\lambda^{-\alpha}$ y logaritmos.
\pasos
\begin{enumerate}[label=\arabic*.]
  \item (a) $10^{-0{,}076}=e^{-0{,}076\ln10}\approx e^{-0{,}175}\approx0{,}8395$.
  \item (b) $\lambda^{-0{,}076}=\frac12$, luego $\lambda=2^{1/0{,}076}=e^{\ln2/0{,}076}\approx e^{9{,}120}\approx9{,}14\cdot10^3$.
  \item (c) $\log L=\alpha\log N_c-\alpha\log N$: pendiente $-0{,}076$.
\end{enumerate}
\resultado \textbf{(a) $\times0{,}8395$ (baja un $16\,\%$); (b) unas $9{,}14\cdot10^3$ veces más; (c) pendiente $-0{,}076$.}
\comprobacion (b) $\log_{10}\lambda=\log_{10}2/0{,}076\approx3{,}961$, es decir, casi cuatro órdenes de magnitud.
\end{solucion}

\begin{solucion}{eje:5-log23}
\datos $\log_23>0$ porque $3>1$.
\planteamiento reducción al absurdo y paridad.
\pasos
\begin{enumerate}[label=\arabic*.]
  \item Si $\log_23=p/q$ con $p,q\in\N$, entonces $2^{p/q}=3$ y, elevando a $q$, $2^p=3^q$.
  \item $2^p$ es par (porque $p\ge1$) y $3^q$ es impar (producto de impares).
  \item Contradicción.
\end{enumerate}
\resultado \textbf{$\log_23$ es irracional.}
\comprobacion $\log_23\approx1{,}585$; las aproximaciones racionales $3/2$ y $8/5$ dan $2^3=8\neq9=3^2$ y $2^8=256\neq243=3^5$.
\end{solucion}

\begin{solucion}{eje:5-temperatura}
\datos $z\in\R^n$, $M=\max_iz_i$, $T>0$.
\planteamiento aplicar la \cref{prop:lse}(ii) al vector $z/T$, cuyo máximo es $M/T$.
\pasos
\begin{enumerate}[label=\arabic*.]
  \item $\frac MT\le\LSE(z/T)\le\frac MT+\ln n$.
  \item Multiplicando por $T>0$: $M\le T\LSE(z/T)\le M+T\ln n$.
  \item Cuando $T\to0^+$, $T\ln n\to0$; por el criterio del sándwich, $T\LSE(z/T)\to M$.
\end{enumerate}
\resultado \textbf{$M\le T\LSE(z/T)\le M+T\ln n$, y el límite es $M$}: con temperatura baja, log-suma-exp se comporta como el máximo y la softmax se concentra en el logit mayor.
\comprobacion $z=(1,0)$, $T=0{,}1$: $0{,}1\ln(e^{10}+1)\approx1{,}0000045$, entre $1$ y $1+0{,}1\ln2\approx1{,}069$.
\end{solucion}

\begin{solucion}{eje:5-lse-media}
\datos $z\in\R^n$; $\bar z=\frac1n\sum z_i$.
\planteamiento aplicar el \cref{teo:amgm} a los números positivos $e^{z_i}$.
\pasos
\begin{enumerate}[label=\arabic*.]
  \item $\frac1n\sum_ie^{z_i}\ge\sqrt[n]{\prod_ie^{z_i}}=e^{\frac1n\sum_iz_i}=e^{\bar z}$.
  \item Multiplicando por $n$ y tomando $\ln$ (creciente): $\LSE(z)\ge\ln n+\bar z$.
  \item Igualdad si y solo si todos los $e^{z_i}$ son iguales, es decir, todos los $z_i$ son iguales.
\end{enumerate}
\resultado \textbf{$\LSE(z)\ge\bar z+\ln n$, con igualdad si y solo si $z$ es constante.}
\comprobacion con $z=(2,1,0,-1)$: $\bar z+\ln4\approx0{,}5+1{,}386=1{,}886\le2{,}440$.
\end{solucion}

\begin{solucion}{eje:5-gibbs}
\datos $p_i,q_i>0$, $\sum p_i=\sum q_i=1$.
\planteamiento aplicar $\ln x\le x-1$ (\cref{cor:lnx}) con $x=q_i/p_i$ y sumar con pesos $p_i$.
\pasos
\begin{enumerate}[label=\arabic*.]
  \item $\ln\frac{q_i}{p_i}\le\frac{q_i}{p_i}-1$; multiplicando por $p_i>0$: $p_i\ln\frac{q_i}{p_i}\le q_i-p_i$.
  \item Sumando en $i$: $\sum_ip_i\ln\frac{q_i}{p_i}\le\sum_iq_i-\sum_ip_i=1-1=0$.
  \item Como $\ln\frac{q_i}{p_i}=\ln q_i-\ln p_i$, la desigualdad es $\sum_ip_i\ln q_i\le\sum_ip_i\ln p_i$; cambiando de signo se obtiene la forma del enunciado.
\end{enumerate}
\resultado \textbf{$-\sum_ip_i\ln q_i\ge-\sum_ip_i\ln p_i$}: la entropía cruzada es al menos la entropía (desigualdad de Gibbs).
\comprobacion $p=(\frac12,\frac12)$, $q=(\frac14,\frac34)$: $-\frac12(\ln\frac14+\ln\frac34)\approx0{,}8370\ge\ln2\approx0{,}6931$.
\end{solucion}

\begin{solucion}{eje:5-e-creciente}
\datos $n\in\N$; los $n+1$ números $a_1=\dots=a_n=1+\frac1n$ y $a_{n+1}=1$.
\planteamiento \cref{teo:amgm} con desigualdad estricta, porque los números no son todos iguales.
\pasos
\begin{enumerate}[label=\arabic*.]
  \item Media aritmética: $\frac{n(1+\frac1n)+1}{n+1}=\frac{n+2}{n+1}=1+\frac1{n+1}$.
  \item Media geométrica: $\bigl((1+\frac1n)^n\cdot1\bigr)^{1/(n+1)}$.
  \item Como $1+\frac1n\neq1$, la desigualdad es estricta: $\bigl(1+\frac1n\bigr)^{n/(n+1)}<1+\frac1{n+1}$.
  \item Elevando a $n+1$ (números positivos, \cref{eje:1-potencia-n}): $\bigl(1+\frac1n\bigr)^n<\bigl(1+\frac1{n+1}\bigr)^{n+1}$.
\end{enumerate}
\resultado \textbf{La sucesión $(1+\frac1n)^n$ es estrictamente creciente}, lo que es un paso para probar que converge al número~$e$.
\comprobacion $n=1,2,3,4$: $2<2{,}25<2{,}370<2{,}441$.
\end{solucion}
% ---------------------------------------------------------------------
\section{Tema 6}

\begin{solucion}{eje:6-a-decimal}
\planteamiento sumar cada cifra por su potencia de la base (\cref{teo:base}).
\pasos
\begin{enumerate}[label=\arabic*.]
  \item $(1011\,0110)_2=128+32+16+4+2=182$.
  \item $\texttt{0x2F}=2\cdot16+15=47$; $\texttt{0xFF}=15\cdot16+15=255$.
  \item $(777)_8=7\cdot64+7\cdot8+7=511$.
\end{enumerate}
\resultado \textbf{$182$, $47$, $255$ y $511$.}
\comprobacion $(1011\,0110)_2$ agrupado en cuartetos es $\texttt{B}\,\texttt{6}$, y $\texttt{0xB6}=11\cdot16+6=182$; $(777)_8=8^3-1=511$, como predice la \cref{prop:cifras}.
\end{solucion}

\begin{solucion}{eje:6-contar}
\planteamiento \cref{prop:cifras} y \cref{teo:c2}.
\pasos
\begin{enumerate}[label=\arabic*.]
  \item (a) $2^8=256$, $2^{16}=65\,536$, $2^{32}=4\,294\,967\,296$.
  \item (b) Sin signo: $[0,255]$. Complemento a dos: $[-2^7,2^7-1]=[-128,127]$.
  \item (c) Cada cifra hexadecimal son 4 bits: $32/4=8$ cifras.
\end{enumerate}
\resultado \textbf{(a) $256$, $65\,536$, $4\,294\,967\,296$; (b) $[0,255]$ y $[-128,127]$; (c) $8$.}
\comprobacion ambos rangos de (b) contienen $256$ enteros.
\end{solucion}

\begin{solucion}{eje:6-vf}
\planteamiento \cref{prop:fraccion-binaria}, \cref{tab:formatos} y \cref{teo:c2}.
\pasos
\begin{enumerate}[label=(\alph*)]
  \item Falsa: $1/10$ tiene denominador no potencia de~2, su desarrollo binario es infinito.
  \item Verdadera: ambos tienen 8 bits de exponente; máximos $\approx3{,}39\cdot10^{38}$ y $3{,}40\cdot10^{38}$.
  \item Verdadera: $T(w)<0$ si y solo si $a_{n-1}=1$.
  \item Verdadera: $0{,}625=5/8=(0{,}101)_2=(1{,}01)_2\cdot2^{-1}$ necesita solo 2 bits de fracción.
\end{enumerate}
\resultado \textbf{(a) falsa; (b) verdadera; (c) verdadera; (d) verdadera.}
\comprobacion (b) $(2-2^{-7})2^{127}\approx3{,}390\cdot10^{38}$ y $(2-2^{-23})2^{127}\approx3{,}403\cdot10^{38}$.
\end{solucion}

\begin{solucion}{eje:6-a-binario}
\planteamiento divisiones sucesivas entre 2 (\cref{ej:conversion}) y agrupación en cuartetos.
\pasos
\begin{enumerate}[label=\arabic*.]
  \item $200=128+64+8=(1100\,1000)_2=\texttt{0xC8}$.
  \item $1000=512+256+128+64+32+8=(11\,1110\,1000)_2=\texttt{0x3E8}$.
  \item $4095=2^{12}-1=(1111\,1111\,1111)_2=\texttt{0xFFF}$.
\end{enumerate}
\resultado $200=\texttt{0xC8}$, $1000=\texttt{0x3E8}$ y $4095=\texttt{0xFFF}$, con las expresiones binarias de los pasos 1--3.
\comprobacion Horner: $12\cdot16+8=200$; $(3\cdot16+14)\cdot16+8=1000$; $(15\cdot16+15)\cdot16+15=4095$.
\end{solucion}

\begin{solucion}{eje:6-fracciones}
\planteamiento multiplicar sucesivamente por 2 y anotar la parte entera.
\pasos
\begin{enumerate}[label=(\alph*)]
  \item $0{,}625\to1{,}25$ (1), $0{,}25\to0{,}5$ (0), $0{,}5\to1$ (1): $(0{,}101)_2$.
  \item $13/16=(1101)_2/2^4=(0{,}1101)_2$.
  \item $1/3\to2/3$ (0), $2/3\to4/3$ (1), $1/3\to\dots$: $(0{,}\overline{01})_2$.
  \item $0{,}1\to0{,}2$ (0), $0{,}4$ (0), $0{,}8$ (0), $1{,}6$ (1), $1{,}2$ (1), $0{,}4$ (0), \dots: $(0{,}0\overline{0011})_2$.
\end{enumerate}
\resultado \textbf{(a) $0{,}101$; (b) $0{,}1101$; (c) $0{,}\overline{01}$; (d) $0{,}0\overline{0011}$.} Son finitos (a) y (b), de denominadores $8$ y $16$; (c) y (d) tienen denominadores $3$ y $10$, que no son potencias de~2.
\comprobacion (c) $\sum_{k\ge1}4^{-k}=\frac{1/4}{1-1/4}=\frac13$.
\end{solucion}

\begin{solucion}{eje:6-c2}
\planteamiento \cref{def:c2} y \cref{teo:c2}(iii): opuesto $=$ invertir y sumar~1.
\pasos
\begin{enumerate}[label=\arabic*.]
  \item (a) $1=(0000\,0001)_2$, invertido $1111\,1110$, más 1: $-1=(1111\,1111)_2$.
  \item $100=(0110\,0100)_2$, invertido $1001\,1011$, más 1: $-100=(1001\,1100)_2$. $127=(0111\,1111)_2$.
  \item (b) $(1000\,0001)_2=-128+1=-127$; $(1111\,0110)_2=-128+64+32+16+4+2=-10$.
  \item (c) $100+50=150>127$: desbordamiento. Se obtiene $150-256=-106$, es decir, $(1001\,0110)_2$.
\end{enumerate}
\resultado \textbf{(a) $1111\,1111$, $1001\,1100$, $0111\,1111$; (b) $-127$ y $-10$; (c) el resultado de 8 bits es $-106$ (desbordamiento).}
\comprobacion $(1001\,1100)_2=-128+16+8+4=-100$; $(1001\,0110)_2=-128+16+4+2=-106$.
\end{solucion}

\begin{solucion}{eje:6-utf8}
\planteamiento escribir el punto de código en binario con 11 o 16 bits y repartir los bits en las plantillas.
\pasos
\begin{enumerate}[label=\arabic*.]
  \item «é»: $\texttt{0xE9}=(000\,1110\,1001)_2$ con 11 bits; se reparte como $00011\mid101001$.
  \item Bytes: $110\,00011=\texttt{0xC3}$ y $10\,101001=\texttt{0xA9}$.
  \item «€»: $\texttt{0x20AC}=(0010\,0000\,1010\,1100)_2$; se reparte como $0010\mid000010\mid101100$.
  \item Bytes: $1110\,0010=\texttt{0xE2}$, $10\,000010=\texttt{0x82}$, $10\,101100=\texttt{0xAC}$.
\end{enumerate}
\resultado «é» se codifica con los bytes \texttt{C3 A9} y «€» con \texttt{E2 82 AC}.
\comprobacion al decodificar \texttt{C3 A9}, los bits útiles $00011$ y $101001$ dan $3\cdot64+41=233=\texttt{0xE9}$.
\end{solucion}

\begin{solucion}{eje:6-digitos}
\planteamiento número de cifras $=\lfloor\log_bn\rfloor+1$ (\cref{prop:cifras}).
\pasos
\begin{enumerate}[label=\arabic*.]
  \item (a) $2^{29}=536\,870\,912\le10^9<2^{30}=1\,073\,741\,824$: $\lfloor\log_210^9\rfloor=29$ y hay $30$ cifras.
  \item (b) $\log_{10}2^{100}=100\cdot0{,}30103=30{,}103$: $31$ cifras.
  \item (c) $2^{32}-1=(\texttt{FFFFFFFF})_{16}$: $8$ cifras.
\end{enumerate}
\resultado \textbf{(a) $30$; (b) $31$; (c) $8$.}
\comprobacion (b) $2^{100}=1\,267\,650\,600\,228\,229\,401\,496\,703\,205\,376$ tiene 31 cifras.
\end{solucion}

\begin{solucion}{eje:6-decodificar}
\planteamiento separar signo, exponente y fracción (\cref{def:flotante}); sesgo $127$ en fp32 y bf16.
\pasos
\begin{enumerate}[label=\arabic*.]
  \item $\texttt{0x3F800000}=0\ 01111111\ 000\dots0$: $s=0$, $E=127$, $f=0$; valor $1{,}0\cdot2^0=1$.
  \item $\texttt{0xC0A00000}=1\ 10000001\ 0100\dots0$: $s=1$, $E=129$, $f=(0{,}01)_2=0{,}25$; valor $-1{,}25\cdot2^{2}=-5$.
  \item $0{,}75=(1{,}1)_2\cdot2^{-1}$: $E=126=(01111110)_2$, $f=1000000$. Palabra $0\,01111110\,1000000=(0011\,1111\,0100\,0000)_2=\texttt{0x3F40}$.
  \item $3=(1{,}1)_2\cdot2^{1}$: $E=128=(10000000)_2$. Palabra $0\,10000000\,1000000=\texttt{0x4040}$.
\end{enumerate}
\resultado (a) $1$ y $-5$; (b) en bf16, $0{,}75$ es \texttt{0x3F40} y $3$ es \texttt{0x4040}.
\comprobacion los 16 bits superiores de los fp32 de $0{,}75$ ($\texttt{0x3F400000}$) y $3$ ($\texttt{0x40400000}$) coinciden con las palabras bf16, como se explicó tras la \cref{tab:formatos}.
\end{solucion}

\begin{solucion}{eje:6-espaciado}
\datos bf16: $p=7$; fp32: $p=23$; fp16: $e=5$, $p=10$, sesgo $15$.
\planteamiento en $[2^k,2^{k+1})$ los números son $(1{,}f)_2\cdot2^k$ con $f$ de $p$ bits.
\pasos
\begin{enumerate}[label=\arabic*.]
  \item (a) Al aumentar $f$ en una unidad del último bit, el número aumenta $2^{-p}\cdot2^k=2^{k-7}$.
  \item (b) En $[256,512)=[2^8,2^9)$ la distancia es $2^{1}=2$: los representables cercanos son $256$ y $258$. $257$ está a la misma distancia; el empate se resuelve hacia la fracción par, $256$ ($f=0000000$). Luego $\fl(256+1)=256$ y $\fl(\fl(256+1)+1)=\fl(257)=256$.
  \item $\fl(1+1)=2$ y $\fl(256+2)=258$, que es representable.
  \item (c) Hasta $2^8=256$ la distancia es $\le1$ y todos los enteros son representables; $257$ es el primero que no lo es. En fp32, por el mismo argumento, $2^{24}+1=16\,777\,217$.
  \item (d) El mayor exponente útil de fp16 es $E=2^5-2=30$, es decir, $2^{30-15}=2^{15}$, con fracción máxima $1-2^{-10}$: $(2-2^{-10})\cdot2^{15}=65\,536-32=65\,504$.
\end{enumerate}
\resultado \textbf{(a) $2^{k-7}$; (b) $256$, $256$ y $258$: la suma no es asociativa; (c) $257$ en bf16 y $16\,777\,217$ en fp32; (d) $65\,504$.}
\comprobacion (c) en bf16, $256=(1{,}0000000)_2\cdot2^8$ y $258=(1{,}0000001)_2\cdot2^8$ son consecutivos.
\end{solucion}

\begin{solucion}{eje:6-memoria}
\datos $7\cdot10^9$ parámetros; 4, 2, 1 y $0{,}5$ bytes por parámetro.
\planteamiento memoria $=$ parámetros $\times$ bytes por parámetro.
\pasos
\begin{enumerate}[label=\arabic*.]
  \item fp32: $2{,}8\cdot10^{10}$ bytes $=28$ GB; en GiB, $2{,}8\cdot10^{10}/2^{30}\approx26{,}08$ GiB.
  \item bf16: $14$ GB. int8: $7$ GB. 4 bits: $3{,}5$ GB.
\end{enumerate}
\resultado \textbf{$28$ GB ($26{,}08$ GiB), $14$ GB, $7$ GB y $3{,}5$ GB.}
\comprobacion cada reducción a la mitad de bits divide la memoria por~2; $28/26{,}08\approx1{,}074$, el factor entre GiB y GB.
\end{solucion}

\begin{solucion}{eje:6-divisibilidad}
\datos $n=\sum_kd_k10^k$.
\planteamiento por el \cref{eje:3-diferencia-potencias}, $10^k-1=(10-1)(10^{k-1}+\dots+1)$ es múltiplo de~9 y, por tanto, de~3.
\pasos
\begin{enumerate}[label=\arabic*.]
  \item $n-\sum_kd_k=\sum_kd_k(10^k-1)$, que es múltiplo de~9.
  \item Si dos enteros difieren en un múltiplo de~9 (o de~3), uno es múltiplo de~9 (o de~3) si y solo si el otro lo es.
  \item Base 16: $16^k-1=(16-1)(16^{k-1}+\dots+1)$ es múltiplo de 15, luego $n-\sum_kh_k$ es múltiplo de 15 si $n=\sum_kh_k16^k$.
  \item Criterio: $n$ es divisible por 15 (y, por tanto, por 3 y por 5) si y solo si lo es la suma de sus cifras hexadecimales.
\end{enumerate}
\resultado \textbf{Los criterios son ciertos en base 10 (para 3 y 9) y en base 16 (para 15, 3 y 5).}
\comprobacion $\texttt{0x3C}=60$: $3+12=15$, y $60=4\cdot15$.
\end{solucion}

\begin{solucion}{eje:6-extension}
\datos $w=(a_{n-1}\dots a_0)$, $w'=(a_{n-1}a_{n-1}a_{n-2}\dots a_0)$ con $n+1$ bits.
\planteamiento aplicar la \cref{def:c2} a $w'$ y separar casos según el bit de signo.
\pasos
\begin{enumerate}[label=\arabic*.]
  \item $T(w')=-a_{n-1}2^{n}+a_{n-1}2^{n-1}+\sum_{k=0}^{n-2}a_k2^k$.
  \item $-2^n+2^{n-1}=-2^{n-1}$, luego $T(w')=-a_{n-1}2^{n-1}+\sum_{k=0}^{n-2}a_k2^k=T(w)$.
\end{enumerate}
\resultado \textbf{$T(w')=T(w)$}; por inducción, la extensión a cualquier número de bits conserva el valor.
\comprobacion $-5=(1011)_2$ en 4 bits y $(1111\,1011)_2$ en 8 bits (\cref{ej:c2}).
\end{solucion}

\begin{solucion}{eje:6-decimal-finito}
\datos $p/q$ irreducible, $q\ge1$.
\planteamiento como en la \cref{prop:fraccion-binaria}, con $10=2\cdot5$ y el teorema fundamental de la aritmética.
\pasos
\begin{enumerate}[label=\arabic*.]
  \item Si el desarrollo tiene $k$ cifras decimales, $p/q=N/10^k$, luego $p\,10^k=qN$ y $q\mid p\,10^k$.
  \item Como $q$ no comparte factores primos con $p$, $q\mid10^k=2^k5^k$, y sus únicos factores primos posibles son 2 y 5: $q=2^a5^b$.
  \item Recíprocamente, si $q=2^a5^b$ y $m=\max(a,b)$, entonces $\frac pq=\frac{p\,2^{m-a}5^{m-b}}{10^m}$, que tiene a lo sumo $m$ cifras decimales tras la coma.
\end{enumerate}
\resultado \textbf{$p/q$ tiene desarrollo decimal finito si y solo si $q=2^a5^b$.}
\comprobacion $\frac3{40}=\frac3{2^35}=\frac{75}{1000}=0{,}075$; $\frac16$ tiene el factor~3 y es periódico.
\end{solucion}

\begin{solucion}{eje:6-desbordamiento}
\datos $x,y\in[-2^{n-1},2^{n-1}-1]$; $s\equiv x+y\pmod{2^n}$ con $s$ en el rango.
\planteamiento acotar $x+y$ en cada caso; hay desbordamiento si y solo si $x+y$ sale del rango.
\pasos
\begin{enumerate}[label=\arabic*.]
  \item (a) Si $x\ge0>y$: $y\le x+y<x$, luego $x+y$ está en el rango y $s=x+y$. El caso $y\ge0>x$ es simétrico.
  \item (b) Si $x,y\ge0$: $0\le x+y\le2^n-2$. Si $x+y\le2^{n-1}-1$, no hay desbordamiento y $s\ge0$. Si $x+y\ge2^{n-1}$, hay desbordamiento y $s=x+y-2^n\in[-2^{n-1},-2]$, negativo.
  \item Si $x,y<0$: $-2^n\le x+y\le-2$. Si $x+y\ge-2^{n-1}$, $s=x+y<0$. Si $x+y<-2^{n-1}$, hay desbordamiento y $s=x+y+2^n\in[0,2^{n-1}-1]$, no negativo.
  \item Juntando los casos: hay desbordamiento exactamente cuando los sumandos tienen el mismo signo y $s$ tiene el contrario.
\end{enumerate}
\resultado \textbf{Sumandos de signo distinto nunca desbordan; el desbordamiento equivale a un cambio de signo con sumandos del mismo signo.}
\comprobacion $n=8$: $100+50$ da $s=-106<0$ (\cref{eje:6-c2}).
\end{solucion}

\begin{solucion}{eje:6-redondeo}
\datos $2^k\le\abs x<2^{k+1}$; $p$ bits de fracción; $\fl$ redondea al más cercano.
\planteamiento los números representables en $[2^k,2^{k+1}]$ forman una rejilla de paso $2^{k-p}$ (\cref{eje:6-espaciado}); el error es como mucho medio paso.
\pasos
\begin{enumerate}[label=\arabic*.]
  \item Sin pérdida de generalidad $x>0$. Los números $2^k+j\,2^{k-p}$, $j=0,\dots,2^p$, son representables (para $j=2^p$ se obtiene $2^{k+1}$).
  \item $x$ está entre dos consecutivos, $c\le x\le c+2^{k-p}$, y el más cercano dista como mucho $\frac12\cdot2^{k-p}=2^{k-p-1}$.
  \item Como $\fl(x)$ es el más cercano de todo el formato, $\abs{\fl(x)-x}\le2^{k-p-1}$.
  \item Dividiendo por $\abs x\ge2^k$: $\frac{\abs{\fl(x)-x}}{\abs x}\le\frac{2^{k-p-1}}{2^k}=2^{-(p+1)}$.
  \item Escribiendo $\delta=(\fl(x)-x)/x$ se obtiene $\fl(x)=x(1+\delta)$ con $\abs\delta\le u$.
\end{enumerate}
\resultado \textbf{$\abs{\fl(x)-x}\le2^{k-p-1}$ y el error relativo es $\le u=2^{-(p+1)}$.}
\comprobacion en bf16, $x=257$: $k=8$, $p=7$; $\abs{256-257}=1=2^{8-7-1}$ y $1/257\approx0{,}003891\le2^{-8}\approx0{,}003906$.
\end{solucion}
